\documentclass[11pt,a4paper]{article}
\usepackage[margin=17mm,top=13mm,bottom=11mm]{geometry}
\usepackage{amsmath,bm,microtype,xcolor,booktabs}
\usepackage[most]{tcolorbox}
\usepackage{listings}
\usepackage{enumitem}
\setlist[itemize]{itemsep=2pt,topsep=2pt,leftmargin=1.3em}
\definecolor{hl}{RGB}{0,80,150}
\definecolor{cbg}{RGB}{246,246,240}
\lstset{basicstyle=\ttfamily\scriptsize,backgroundcolor=\color{cbg},frame=single,
        rulecolor=\color{black!25},breaklines=true,columns=fullflexible,
        xleftmargin=2pt,xrightmargin=2pt,framesep=4pt,aboveskip=4pt,belowskip=4pt}
\newcommand{\code}[1]{\texttt{\small #1}}
\newcommand{\Told}{T^{\rm old}}\newcommand{\Tnew}{T^{\rm new}}
\pagestyle{empty}
\begin{document}

{\large\bfseries The zero-sum sheath stabiliser in JOREK}\\[1pt]
{\footnotesize where it lives, what it looks like, and what you set ---
after SOLPS-ITER \code{b2stbc\_stab\_coeff\_sheath\_te/ti/ni}}
\vspace{3pt}\hrule\vspace{7pt}

\textbf{1. The input parameters.} Three reals in namelist \code{in1}, all default \textbf{0},
which reproduces the previous code bit-for-bit:
\begin{lstlisting}
stab_coeff_sheath_te = 100.d0   ! electron energy sheath row
stab_coeff_sheath_ti = 100.d0   ! ion  energy sheath row
stab_coeff_sheath_ni = 100.d0   ! density     sheath row
\end{lstlisting}
\vspace{-3pt}
{\small
\begin{tabular}{@{}ll@{}}
declared & \code{models/phys\_module.f90}\\
defaults & \code{models/preset\_parameters.f90:81--83}\\
namelist & \code{models/model600/initialise\_parameters.f90:67--68}\\
broadcast & \code{communication/broadcast\_phys.f90} (packed/unpacked in matching order)\\
\end{tabular}}

\vspace{7pt}
\textbf{2. The maths.} Let $\Pi$ be the sheath \emph{collection flux} --- the geometric factor
already multiplying every term in that boundary row:
\[
\Pi \;=\;\underbrace{(\bm B\!\cdot\!\bm n)\,V_\parallel R\,\mathrm{d}l}_{\text{parallel}}
\;+\;\underbrace{v_{E,n}^{\rm clip}\,R\,\mathrm{d}l}_{E\times B}
\;+\;\underbrace{c_s R\,\mathrm{d}l\,c_\alpha}_{\text{grazing floor}}
\]
The \emph{physical} sheath energy source and the \emph{added} stabiliser are
\[
S_{T_i}= -(\gamma_{\rm sh,i}\!-\!1)\,\rho\,T_i\,\Pi ,
\qquad\qquad
\boxed{\,S^{\rm stab}_{T_i}= +\,\alpha_i\,\rho\,\Pi\,\bigl(\Told_i-\Tnew_i\bigr)\,}
\]
JOREK linearises about the old state, so with $b$ the residual and $\mathsf A$ the matrix:
\[
S^{\rm stab}\big|_{\Tnew=\Told}=0
\;\Rightarrow\;\textbf{no residual contribution},
\qquad
\frac{\partial S^{\rm stab}}{\partial \Tnew}=-\alpha_i\rho\Pi
\;\Rightarrow\;
\mathsf A_{T_iT_i}\mathrel{+}=\theta\,\alpha_i\,\rho\,\Pi\,\hat T_i
\]
so the assembled column becomes $(\gamma_{\rm sh,i}-1+\alpha_i)\rho\hat T_i\Pi$ while the residual
keeps $(\gamma_{\rm sh,i}-1)$. \textcolor{hl}{\textbf{That asymmetry is the entire
implementation.}} The converged state is where $b=0$, and $\mathsf A$ does not enter that
condition --- so $\alpha$ cannot move the answer.

\vspace{7pt}
\textbf{3. The code.} One file: \code{models/model600/mod\_boundary\_matrix\_open.f90}.

\vspace{2pt}
{\small\textit{The residual --- untouched, keeps the physical $\gamma$:}}
\begin{lstlisting}
rhs_ij(var_Ti) = - v * (gamma_sheath_i-1.d0) * r0 * Ti0 * vpar0 * ps0_s * normal_sign3 * tstep &
                 - v * (gamma_sheath_i-1.d0) * r0 * Ti0 * fu_ven_sh * BigR * dl        * tstep &
                 - v * (gamma_sheath_i-1.d0) * r0 * Ti0 * cs0    * BigR * dl * c_angle * tstep &
\end{lstlisting}
{\small\textit{The Jacobian --- the added term (line 651). The bracket is $\Pi$: the same three
factors as above.}}
\begin{lstlisting}
amat(var_Ti,var_Ti) = amat(var_Ti,var_Ti)                                    &
                    + v * stab_coeff_sheath_ti * r0 * Ti                     &
                      * ( vpar0 * ps0_s * normal_sign3                       &  ! parallel
                        + fu_ven_sh * BigR * dl                              &  ! clipped ExB
                        + cs0 * BigR * dl * c_angle )                        &  ! grazing floor
                      * theta * tstep
\end{lstlisting}
\vspace{-2pt}
{\small \code{v} = test function, \code{Ti} = \emph{trial} function, \code{r0}/\code{vpar0}/%
\code{cs0} = lagged old-state values, \code{theta*tstep} = the standard implicitness weighting
used by every neighbouring entry.}

\vspace{5pt}
\textbf{4. All four insertion points} --- Jacobian diagonal only, nothing off-diagonal, no other
equation:
\begin{center}\small
\begin{tabular}{clll}
\toprule
line & matrix entry & parameter & row \\
\midrule
616 & \code{amat(var\_rho,var\_rho)} & \code{stab\_coeff\_sheath\_ni} & density \\
651 & \code{amat(var\_Ti,var\_Ti)}   & \code{stab\_coeff\_sheath\_ti} & ion energy \\
673 & \code{amat(var\_Te,var\_Te)}   & \code{stab\_coeff\_sheath\_te} & electron energy \\
706 & \code{amat(var\_T,var\_T)}     & \code{stab\_coeff\_sheath\_ti} & single-$T$ build \\
\bottomrule
\end{tabular}
\end{center}

\vspace{3pt}
\textcolor{red!60!black}{\textbf{Careful:}} had we scaled \emph{both} residual and Jacobian we
would simply have raised $\gamma_{\rm sheath}$ --- a different sheath and a different steady state.
Only the Jacobian is touched.

\vspace{6pt}
\textbf{5. What it does, in numbers.} Same equation
$\bigl(\Tnew-\Told\bigr)/\Delta t=Q-kT+\alpha k(\Told-\Tnew)$ with $Q=5$, $k=2$,
$\Delta t=0.5$, start $T=9$, so $T^{*}=Q/k=2.5$:
\begin{center}\small
\begin{tabular}{lccccccc c}
\toprule
& start & 1 & 2 & 3 & 4 & 5 & 6 & after 4000 steps\\
\midrule
$\alpha=0$   & 9.000 & 5.750 & 4.125 & 3.313 & 2.906 & 2.703 & 2.602 & \textbf{2.5000000000}\\
$\alpha=100$ & 9.000 & 8.936 & 8.873 & 8.811 & 8.749 & 8.688 & 8.627 & \textbf{2.5000000000}\\
\bottomrule
\end{tabular}
\end{center}
\vspace{-3pt}
Identical destination to ten digits; $\alpha$ only takes smaller bites, so no single step can
overshoot.

\vspace{5pt}
\textbf{6. Verification.} \code{test\_sheath\_stabiliser} in
\code{tests/floating\_transport/test\_boundary.f90} assembles at $\alpha=0,10,20$ and asserts:
residual \textbf{bit-identical}; Jacobian \textbf{exactly linear} in $\alpha$; \textbf{nothing}
outside the $T_i$/$T_e$/$\rho$ diagonal blocks changes.

\vspace{5pt}
\textbf{7. Result.} Drift-enabled run: previously terminated at $\le649$ steps
$\;\rightarrow\;$ \textbf{past 1100 and converging}.

\vspace{3pt}
{\footnotesize\textcolor{black!55}{Note: SOLPS iterates within a time step, so there the term
vanishes before the step is accepted and is purely a convergence aid. JOREK takes one solve per
step, so $\Told$ is the previous \emph{time step} and genuine transients are slowed as well. The
stationary state is unaffected either way.}}

\end{document}
